\documentclass[11pt,a4paper]{article}
\usepackage[top=42pt,bottom=42pt,left=70pt,right=70pt]{geometry}  % to make pages wider

\usepackage{amsmath}
\usepackage{amsfonts}
\usepackage{graphicx}  % for figures etc.
\usepackage{url}  % to handle URLs 
\usepackage[comma,authoryear]{natbib}  % for author-date ref style --  allows \citet (textual), and \citep (parenthesized)
\usepackage[hidelinks]{hyperref}
\hypersetup{allcolors=[rgb]{0, 0, 0.33},colorlinks=true}

\renewcommand {\cite} {\citep}  % default for cite is citet in natbib - changes it to citep
\renewcommand\bibname{References}    % if not using newrucsthesis sty file

\usepackage{booktabs}
\usepackage{array}
\usepackage{enumitem}
\usepackage{listings}
\usepackage{amssymb}

\usepackage{setspace}
\doublespacing

% Optional: Customizing how the code looks
\lstset{
    basicstyle=\ttfamily\small,
    numbers=left,           
    numberstyle=\tiny,     
    stepnumber=2,          
    numbersep=5pt,
    frame=single,           
    rulecolor=\color{black},
    frameround=tttt,        
    breaklines=true         
}

\title{Implementation}
\author{g23n7880 }
\date{September 2026}

\begin{document}


\section{Overview}
This chapter documents the concrete realisation of the semi-automated pipeline described in Chapter 3, presenting the implementation of each stage together. The pipeline was implemented in Python 3.11 across a set of modular scripts, each corresponding to a discrete stage. The stages are semi-automated data collection, language identification, quality filtering, and orthographic classification.



\section{Stage 1: Semi-automated Data Collection}

Data collection used a semi-automated approach, with source selection performed manually by the researcher. At the same time, all subsequent processing, including fetching, extraction, filtering, and labelling, was handled by the pipeline. 

Sources were configured as structured tuples containing the URL, the orthographic variant label, a source identifier, an optional CSS selector for targeted extraction, and an optional list of skip patterns for excluding boilerplate content. The following shows the SAS and LS static source configurations:

\begin{lstlisting}[language=Python, caption={Seed source configuration for SAS and LS variants.}, label={lst:sources}]
[language=Python, caption={Seed source configuration for SAS and LS variants.}, label={lst:sources}]
SAS_STATIC_SOURCES = [
    ("https://www.gov.za/st",              "SAS", "govza_sesotho",  None, None),
    ("https://www.sabc.co.za/sabc/sesotho/","SAS","sabc_sesotho",   None, None),
    ("https://www.vukuzenzele.gov.za",     "SAS", "vukuzenzele",    None, None),
    ("https://www.dbe.gov.za/sesotho",     "SAS", "dbe_sesotho",    None, None),
    ("https://www.saps.gov.za/sesotho",    "SAS", "saps_sesotho",   None, None),
]
 
LS_STATIC_SOURCES = [
    ("https://www.gov.ls",                 "LS",  "gov_ls",         None, None),
    ("https://www.lestimes.com",           "LS",  "lestimes",       None, None),
    ("https://www.thepost.co.ls",          "LS",  "thepost_ls",     None, None),
    ("https://www.leihlolabasotho.co.ls",  "LS",  "leihlolabasotho",None, None),
]
\end{lstlisting}

Beyond the static source lists, three additional source-specific collection methods were implemented: a systematic scraper for the Nal'ibali reading campaign, a downloader for the Vuk'uzenzele DSFI GitHub corpus, and a local PDF processor for Sesotho matric examination papers.

\subsection{Web extraction}

HTML text was extracted using the \texttt{fetch\_and\_collect} function, which handles URL fetching, content selection, paragraph extraction, and the quality and LID filtering chain in a single call. The function accepts an optional CSS class selector for cases where the target content is contained within a specific page element, as is the case with the Nalibali website, which renders story content inside \texttt{div} elements carrying the class \texttt{elementor-widget-text-editor}:

\begin{lstlisting}[language=Python, caption={Core web extraction function.}, label={lst:fetch}]
def fetch_and_collect(url, label, source_name, selector=None, skip_patterns=None):
    r = requests.get(url, headers=HEADERS, timeout=15)
    soup = BeautifulSoup(r.text, "html.parser")
 
    if selector:
        divs = soup.find_all("div", class_=selector)
        paragraphs = [p.get_text() for div in divs
                      for p in div.find_all("p")]
    else:
        paragraphs = [p.get_text() for p in soup.find_all("p")]
 
    if skip_patterns:
        paragraphs = [p for p in paragraphs
                      if not any(sp in p.lower()
                                 for sp in skip_patterns)]
 
    raw_text = " ".join(paragraphs)
    sentences = split_sentences(clean_text(raw_text))
 
    results = []
    for sentence in sentences:
        if not is_quality(sentence):
            continue
        ok, conf = verify_sesotho(sentence)
        if ok:
            results.append({"text": sentence, "label": label,
                             "lid_confidence": conf,
                             "source": source_name,
                             "source_url": url})
    return results
\end{lstlisting}

The Nalibali scraper first traversed the eight category pages to collect story URLs dynamically, then called \texttt{fetch\_and\_collect} on each story using the \texttt{elementor-widget-text-editor} selector and a list of skip patterns to exclude author attributions, translator credits, and comprehension question paragraphs. This approach reflects the structural reality of modern web publishing: content is rarely found in raw \texttt{<p>} tags at the page root. However, it is embedded within component frameworks that require targeted selection for clean extraction.

\subsection{PDF extraction}

Government documents fetched from remote URLs using \texttt{requests} library, and Sesotho matric examination papers (National Senior Certificate, Grade 12 Home Language and First Additional) loaded from local storage. PyMuPDF (\texttt{fitz}) was used for text extraction in both cases. 

Examination papers required additional handling to suppress non-Sesotho content. An english-density filter was applied at the page level where pages with more that 15\% of words matched the English marker lexicon were skipped entirely, retaining only the pages dominated by Sesotho content:

\begin{lstlisting}[language=Python, caption={English-density page filter for exam PDF extraction.}, label={lst:pdf}]
for page_num in range(2, total_pages - 1):
    page_text = doc[page_num].get_text("text")
    eng_count = len([
        w for w in re.findall(r'\b[a-zA-Z]{4,}\b', page_text.lower())
        if w in ENGLISH_MARKERS
    ])
    total_words = len(page_text.split())
    if total_words > 0 and eng_count / total_words > 0.15:
        continue  # skip English-heavy pages
    # process remaining pages normally
\end{lstlisting}

The first two pages of each PDF were unconditionally skipped to exclude cover pages and rubrics.

\subsection{Vuk'uzenzele DSFI Corpus}

The Vuk'uzenzele corpus was accessed directly from the DSFI Github repository via the GitHub Contents API, which returns a JSON listing of files in the \texttt{data/simple\_align\_output} directory. Each CSV file in this directory contains sentence-aligned English--Sesotho pairs. The Sesotho column was identified by searching for column names containing the substrings \texttt{sot}, \texttt{sesotho}, or \texttt{target}, with a fallback to the second colum for files where no match was found:

\begin{lstlisting}[language=Python, caption={Vuk'uzenzele Sesotho column extraction.}, label={lst:vuku}]
sesotho_col = None
for col in df_raw.columns:
    if any(x in col.lower() for x in ["sot", "sesotho", "target"]):
        sesotho_col = col
        break
if sesotho_col is None and len(df_raw.columns) >= 2:
    sesotho_col = df_raw.columns[1]
\end{lstlisting}


\subsection{Collection Results}

\begin{table}[ht]
    \centering
    \label{tab:corpus_sources}
    \caption{Corpus Source breakdown after all quality filters.}
    \begin{tabular}{llrr}
     \toprule
    \textbf{Source} & \textbf{Label} & \textbf{Domain} & \textbf{Sentences} \\
    \midrule
    Matric HL Exam (NSC)   & SAS & Educational  & 3,347 \\
    Nalibali               & SAS & Narrative    & 3,165 \\
    Vuk'uzenzele DSFSI     & SAS & Government   & 1,182 \\
    gov.za PDFs            & SAS & Government   &   833 \\
    \midrule
    \textit{SAS subtotal}  &     &              & \textit{5,362} \\
    \midrule
    Mokhosi SN (Zenodo)    & LS  & News         &    99 \\
    Lesotho web sources    & LS  & Mixed        & 3,165 \\
    \midrule
    \textit{LS subtotal}   &     &              & \textit{3,264} \\
    \midrule
    \textbf{Total}         &     &              & \textbf{8,626} \\
    \bottomrule
    \end{tabular}
\end{table}



\section{Stage 1: Language Identification and Quality Filtering}

\subsection{Text cleaning}

Raw text from all sources underwent a four-step cleaning procedure before LID verification. Unicode normalisation to NFC form was applied first, ensuring that diacritical marks are represented in their canonical composed form rather than as decomposed sequences, which is critical for LS content, where diacritics encode phonological distinctions that SAS orthography omits \citep{matlosa2017sesotho}. URL strings were then removed, followed by HTML entity decoding and whitespace normalisation. Question numbers that appear at the start of sentences, a common artefact of PDF extraction of examination papers, were stripped using a regular expression.

\begin{lstlisting}[language=Python, caption={Text cleaning function.}, label={lst:clean}]
def clean_text(text):
    text = unicodedata.normalize("NFC", str(text))
    text = re.sub(r"http\S+", "", text)
    text = re.sub(r"&[a-z]+;", "", text)
    # strip leading question numbers e.g. "3." "(a)" "[1]"
    text = re.sub(r"^\s*[\d]+[\.\)]\s*", "", text)
    text = re.sub(r"^\s*[\(\[]\s*[a-zA-Z\d]\s*[\)\]]\s*", "", text)
    text = re.sub(r"\s+", " ", text).strip()
    return text
\end{lstlisting}

\subsection{Code-Switching Detection}

A code-switching filter was applied to reject sentences containing substantial English content. The filter counts occurrences of unambiguously English words of three or more characters and rejects sentences where this count exceeds two. The three-character minimum is due to single-character tokens such as \textit{a}, \textit{o}, and \textit{e}, which are valid Sesotho morphemes serving as verb agreement prefixes and conjunctions, and would generate false positives under a naive English word count:

\begin{lstlisting}[language=Python, caption={Code-switching detection function.}, label={lst:codeswitching}]
ENGLISH_MARKERS = {
    'the', 'and', 'for', 'are', 'but', 'not', 'you', 'all',
    'can', 'will', 'that', 'this', 'they', 'from', 'have',
    # ... (full set of 60+ unambiguous English function words)
}
 
def is_code_switched(text, max_english=2):
    words = re.findall(r'\b[a-zA-Z]{3,}\b', text.lower())
    hits = [w for w in words if w in ENGLISH_MARKERS]
    return len(hits) > max_english
\end{lstlisting}

An additional filter rejected sentences containing \LaTeX\ anonymisation placeholders of the form \texttt{\$T\$} or \texttt{\$X\$}, which were present in a subset of the Mokhosi SN dataset where named entities needed to be anonymised for sentiment annotation. Sentences where fewer than 60\% of characters were alphabetic were also rejected to exclude fragments consisting predominantly of digits, punctuation or formatting symbols.


\subsection{GlotLID Language Verification}



\bibliographystyle{ruauthordate}
\bibliography{implementation}  
\end{document}
